Threat actors can be organized by skill and resource levels, ranging from low or moderately skilled and resourced actors who may be able to build a malicious product to `advanced persistent threats' (APTs): highly skilled and well-resourced (e.g. state-sponsored) groups with long-term agendas \cite{solaiman2019release}.

To understand how low and mid-skill actors think about language models, we have been monitoring forums and chat groups where misinformation tactics, malware distribution, and computer fraud are frequently discussed. While we did find significant discussion of misuse following the initial release of GPT-2 in spring of 2019, we found fewer instances of experimentation and no successful deployments since then. Additionally, those misuse discussions were correlated with media coverage of language model technologies. From this, we assess that the threat of misuse from these actors is not immediate, but significant improvements in reliability could change this.

Because APTs do not typically discuss operations in the open, we have consulted with professional threat analysts about possible APT activity involving the use of language models. Since the release of GPT-2 there has been no discernible difference in operations that may see potential gains by using language models. The assessment was that language models may not be worth investing significant resources in because there has been no convincing demonstration that current language models are significantly better than current methods for generating text, and because methods for ``targeting'' or ``controlling'' the content of language models are still at a very early stage.


